
\documentclass[a4paper,12pt]{letter}
\usepackage[utf8]{inputenc}
\usepackage[english,slovene]{babel}
%\usepackage{blindtext}
%\usepackage{pdflscape}
\usepackage[tikz]{bclogo}
\usepackage{qrcode,marvosym}
\usepackage{pgf,tikz,pgfplots}
\pgfplotsset{compat=1.14}
\usepackage{mathrsfs}
\usetikzlibrary{arrows}
\usepackage{graphicx}
\newcommand{\ds}{\displaystyle}
\usepackage{fancybox}
\usepackage{amsfonts}
\usepackage{amsmath}
%\usepackage{color}
\usepackage{wallpaper}
\usepackage{hyperref}
\usepackage[cp1250]{inputenc}
\usepackage{array}%%%%%%%%%%%%%%%%%%%%%%%%%%%%%%%%%%%%%%%%%%%%
\newcolumntype{M}[1]{>{\centering\arraybackslash}m{#1}}%%%%
\newcolumntype{N}{@{}m{0pt}@{}}%%%%%%%%%%%%%%%%%%%%%%%%%%%%
   \usepackage{fancyhdr,enumerate}

 \newcommand{\ora}{\overrightarrow}
 \renewcommand{\headrulewidth}{3pt}
 \renewcommand{\headrule}{\hbox to\headwidth{%
   \color{gray}\leaders\hrule height \headrulewidth\hfill}}
    \renewcommand{\footrulewidth}{2pt}
    \renewcommand{\footrule}{\hbox to\headwidth{%
  \color{brown}\leaders\hrule height \footrulewidth\hfill}}
\usepackage{gensymb}
\usepackage{amssymb} 

  \usepackage{fancyhdr}
  \textwidth 20 true cm
  \oddsidemargin -2 true cm
  \evensidemargin -2 true cm
  \topmargin -3.5 true cm
  \textheight 27.5 true cm
  \linespread{1.2}
    \renewcommand{\headrulewidth}{0.3pt}
   \renewcommand{\footrulewidth}{1.9pt}
  \definecolor{qqwuqq}{rgb}{1,0,0 }
  \definecolor{qqqqff}{rgb}{0,0,1}
  \definecolor{qqffqq}{rgb}{0,1,0}
  \definecolor{ffqqqq}{rgb}{1,0,0}
\definecolor{zzuxux}{rgb}{0,0,0}%{0.92,0.03,0.03}
\usepackage{mdframed}
\usepackage{multicol}
 \definecolor{bgblue}{RGB}{0,0,253}
\usepackage{xcolor}
\usepackage{eurosym}
%%%%%%%%%%%%%%%%%%%%%%%%%%%%%%%%%%%%%%%%%%%%%%%%%%%%%%%%
\newcounter{tocke}
\setcounter{tocke}{0}
\usepackage{tikz-3dplot}
%%%%%%%%%%%%%%%%%%%%%%%%%%%%%%%%%%%%%%%%%%%%%%%%%%%%%%%%%%
\mdfdefinestyle{theoremstyle}{%
linecolor=brown!40,linewidth=1pt,%
frametitlerule=true,%
frametitlebackgroundcolor=brown!50,
innertopmargin=\topskip,
}

\mdfdefinestyle{theoremstyle1}{%
linecolor=brown,middlelinewidth=1pt,%
frametitlerule=true,%
apptotikzsetting={\tikzset{mdfframetitlebackground/.append style={%
shade,left color=black!40, right color=gray!10},
mdfbackground/.append style={
top color=white!100!white,
bottom color=black!5
},
}}, 
frametitlerulecolor=gray,
frametitlerulewidth=2pt,
innertopmargin=\topskip,
innerbottommargin=\topskip,
}
\mdtheorem[style=theoremstyle1]{naloga}{Naloga}
\mdtheorem[style=theoremstyle]{resitev}{Analiza Naloge}
\mdtheorem[style=theoremstyle]{kriterij}{\v{S}tevilo dose\v{z}enih to\v{c}k na testu}
%%%%%%%%%%%%%%%%%%%%%%%%%%%%%%%%%%%%%%%%%%%%%%%%%%%%%%%%%%
\usepackage[table]{xcolor}
\usepackage{titlesec}
\usepackage{venndiagram}
\usepackage[printwatermark]{xwatermark}
\usepackage{xcolor}
\newwatermark[allpages,color=brown!2,angle=45,scale=5,
xpos=-5,ypos=0]{matej.info}


%%%%%%%%%%%%%%%%%%%%%%%%%%%%%%%%%%%%%%%%%%%%%%%%%%%%%%%%%%
\titleformat{\subsection}
  {\normalfont\Large\bfseries\color{brown}}
  {\thesubsection}
  {1em}
  {}

\titleformat{\subsubsection}
  {\normalfont\large\bfseries\color{brown}}
  { \thesubsubsection}
  {1em}
  {}
%%%%%%%%%%%%%%%%%%%%%%%%%%%%%%%%%%%%%%%%%%%%%%%%%%%%%%%%%%
\begin{document}
\pagestyle{fancy}
\renewcommand\headrulewidth{0pt}
\lhead{}\chead{}\rhead{}
\rfoot{\vspace*{-2\baselineskip}}


\renewcommand{\vec}{\overrightarrow}

%%%%%%%%%%%%%%%%%%%%%%55\Large

\everymath{\displaystyle}
\begin{minipage}{20cm}
\begin{bclogo}[couleur=brown!40,
  arrondi=0,ombre=true, couleurOmbre=gray!50,
logo =\bcplume,
  noborder=true]
{\textrm{{\Large{Test 2.2} \Huge{}:
%
%
%%%%%%%%%%%%%%%%%%%%%%%%%%%%%%%%%%%%%%%%%
%%%%%%%%%%%%%%%%%%%%%%%%%%%%%%%%%%%%%%%%%
% TUKAJ VNESI VSEBINO TESTA
\large{Eksponentna in logaritemska funkcija. Zaporedja.} \hfill $PTI-2$}}}
%%%%%%%%%%%%%%%%%%%%%%%%%%%%%
\end{bclogo}
\end{minipage}
\begin{minipage}{20cm}
\begin{bclogo}[couleur=brown!40,ombre=true, couleurOmbre=gray!50,
  arrondi=0,
logo =\bcplume,
  noborder=true]
 \setlength{\parskip}{}{}
{\textsc{{\large{Ime in priimek:}
%%%%%%%%%%%%%%%%%%%%%%%%%%%%%%%%%%%%%%%%%
% TUKAJ VNESI VSEBINO TESTA
\vspace{0.2cm} \rule{15cm}{0.2mm}\hfill }}}
%%%%%%%%%%%%%%%%%%%%%%%%%%%%%%%%%%%%%%%%%
\end{bclogo}
\end{minipage}


\def\tocke2{3+3+3}
\begin{naloga}[\hfill \quad $\tocke2$
\qquad \rightsquigarrow \vert a.\qquad|b.\qquad|c.\qquad|] \addtocounter{tocke}{\tocke2}
Reši enačbo:
\vspace{-0.5cm}
\begin{multicols} {3}
   a) $\log_2(3x-1)=4$

b) $6^x=7^x$

c) $5^x=45$ 
\end{multicols}

\end{naloga}






\newpage
\def\tocke18{5+1}
\begin{naloga}[\hfill \quad $\tocke18$
\qquad \rightsquigarrow \vert a.\qquad\vert b.\qquad|
c.\qquad|] \addtocounter{tocke}{\tocke18}

a) Nariši graf funkcije $f(x)=2^{x}-4$, tako da najprej izračunaš začetno vrednost in ničlo ter določiš asimptoto.

b) Izračunaj vrednost točke $A(3,y)$, če leži na grafu ter jo označi.
\end{naloga}

\begin{tikzpicture}[line cap=round,line join=round,>=triangle 45,x=1.0cm,y=1.0cm]
\begin{axis}[
x=1.0cm,y=1.0cm,
axis lines=middle,
ymajorgrids=true,
xmajorgrids=true,
xmin=-6.5600000000000005,
xmax=6.640000000000001,
ymin=-3.639999999999998,
ymax=6.440000000000002,
xtick={-6.0,-5.0,...,6.0},
ytick={-3.0,-2.0,...,6.0},]
\clip(-6.56,-3.64) rectangle (6.64,6.44);
\end{axis}
\end{tikzpicture}


\newpage
\def\tocke6{4+4}
\begin{naloga}[\hfill \quad $\tocke6$
\qquad \rightsquigarrow \vert a.\qquad| b.\qquad| c.\qquad| ] \addtocounter{tocke}{\tocke6}
V aritmetičnem zaporedju je prvi člen enak 9, četrti člen pa -22. 

a) Izračunaj diferenco zaporedja in tretji člen.

b) Izračunaj vsoto prvih 10 členov tega zaporedja.
\end{naloga}

\vspace{10cm}
\def\tocke6{4+3}
\begin{naloga}[\hfill \quad $\tocke6$
\qquad \rightsquigarrow \vert a.\qquad| b.\qquad| c.\qquad| ] \addtocounter{tocke}{\tocke6}
a) Naj bo $\log a=10, \log b=-1,\log c=8.$ Izračunaj $\log x,$ če je $x=\dfrac{a^2b}{\sqrt[8]{c}}.$

b) Izračunaj: $\log_2 8-\log_3 \sqrt{3}-\ln 1+\log 120-\log 12$
\end{naloga}
 



\newpage
\def\tocke3{3+4+3}
\begin{naloga}[\hfill \quad $\tocke3$
\qquad \rightsquigarrow \vert a.\qquad\vert b.\qquad] \addtocounter{tocke}{\tocke3}
a) V aritmeričnem zaporedju $\rule{1cm}{0.1mm},\;6\;,12\;,\rule{1cm}{0.1mm}\;\rule{1cm}{0.1mm}, \ldots$ dopolni tri manjkajoče člene.

b) Izračunaj 7. člen in vsoto prvih 10 členov.

c) Na katerem mestu je člen $3072\;?$
\end{naloga}



\vfill
\mdtheorem[style=theoremstyle1]{kriterij}{Kriterij ocenjevanja}
\begin{kriterij*}[\hfill \v{s}tevilo vseh to\v{c}k na testu: $\arabic{tocke}$]
 $$
 \begin{array}{||c|c|c|c|c|c||
>{}c|c|}
 \hline\hline
 \textup{ocena}& 1&2&3&4&5&\textrm{uspe\v{s}nost v } \%& \textrm{\bf OCENA}\\
  \hline\hline
 \% & [0,45) &[45,60)&[60,75
 ) &[75,90)    &    [90,100]& \mbox{\phantom{ \huge{$\frac{5}{445}$
 15}}
 } 
 & \\
 \hline
 \end{array} \hspace*{0.5cm}\vspace*{-0.3cm} \qrcode[hyperlink,height=0.8in]{http://www.matej.info}
$$
\vspace*{0.1cm}
\end{kriterij*}

\newpage



\subsection{Rešitve:}
\subsubsection{1. naloga}
a) $\log_2(3x-1)=4$

\[
3x-1=2^4=16 \Rightarrow 3x=17 \Rightarrow x=\frac{17}{3}
\]

b) $6^x=7^x$

\[
\left(\frac{6}{7}\right)^x=1 \Rightarrow x=0
\]

c) $5^x=45$

\[
x=\log_5 45=\frac{\log 45}{\log 5}
\]


\subsubsection{2. naloga}

a) Dana je funkcija $f(x)=2^x-4$

\[
f(0)=1-4=-3
\]

Za ničlo velja:
\[
2^x-4=0 \Rightarrow 2^x=4 \Rightarrow x=2
\]

Horizontalna asimptota:
\[
y=-4
\]

b) Točka $A(3,y)$:

\[
y=f(3)=2^3-4=8-4=4
\]

\[
A(3,4)
\]

\begin{tikzpicture}
\begin{axis}[
    axis lines=middle,
    xlabel={$x$},
    ylabel={$y$},
    xmin=-3, xmax=5,
    ymin=-5, ymax=6,
    grid=both,
    width=10cm,
    height=7cm,
    samples=200
]

% funkcija
\addplot[blue, thick] {2^x - 4};

% horizontalna asimptota
\addplot[dashed, red, thick] {-4};
\node[right, red] at (axis cs:4,-4) {$y=-4$};

% točka A(3,4)
\addplot[only marks, mark=*, blue] coordinates {(3,4)};
\node[right] at (axis cs:3,4) {$A(3,4)$};

% začetna vrednost
\addplot[only marks, mark=*] coordinates {(0,-3)};
\node[left] at (axis cs:0,-3) {$(0,-3)$};

% ničla
\addplot[only marks, mark=*] coordinates {(2,0)};
\node[above] at (axis cs:2,0) {$(2,0)$};

\end{axis}
\end{tikzpicture}
```
\subsubsection{3.naloga }

Zapišemo podatke:  
\( a_1 = 9,\ a_4 = -22 \).  


a)  
Splošna formula za \(n\)-ti člen aritmetičnega zaporedja je  
\[
a_n = a_1 + (n-1)d,
\]  
kjer je \(d\) diferenca.

Za \(n=4\):  
\[
a_4 = a_1 + 3d
\]
\[
-22 = 9 + 3d
\]
\[
3d = -31 \implies d = -\frac{31}{3}.
\]

Tretji člen:  
\[
a_3 = a_1 + 2d = 9 + 2\left(-\frac{31}{3}\right) = 9 - \frac{62}{3} = \frac{27}{3} - \frac{62}{3} = -\frac{35}{3}.
\]

\[
\boxed{d = -\frac{31}{3},\quad a_3 = -\frac{35}{3}}
\]

b) 
Formula za vsoto prvih \(n\) členov:  
\[
S_n = \frac{n}{2} \big( 2a_1 + (n-1)d \big).
\]
Za \(n=10\):  
\[
S_{10} = \frac{10}{2} \left[ 2\cdot 9 + 9\cdot \left(-\frac{31}{3}\right) \right]
= 5 \left[ 18 - 93 \right]
= 5 \times (-75) = -375.
\]

\[
\boxed{S_{10} = -375}
\]



\subsubsection{4.naloga}
a) 
\[
x=\frac{a^2b}{\sqrt[8]{c}}
\Rightarrow \log x=2\log a+\log b-\frac{1}{8}\log c
\]

\[
\log x=2\cdot10+(-1)-\frac{1}{8}\cdot8=18
\]

b)
\[
\log_2 8=3,\quad \log_3\sqrt{3}=\frac12,\quad \ln 1=0
\]

\[
\log 120-\log 12=\log 10=1
\]

\[
3-\frac12+1=\frac{7}{2}
\]

---

\subsubsection{5. naloga}

Zaporedje je aritmetično: $6,12$

\[
d=6
\]

Manjkajoči členi:
\[
0,\ 6,\ 12,\ 18,\ 24
\]

b) 
\[
a_7=a_1+6d=0+36=36
\]

\[
S_{10}=\frac{10}{2}(0+54)=270
\]

c)
\[
a_n=6(n-1)=3072
\Rightarrow n-1=512
\Rightarrow n=513
\]


\end{document}